
\section{Discussion}

The results admit a coherent interpretation centred on three findings: the failure of the domestic Phillips curve, the dominance of the ERM\,II spillover channel in-sample, and the inability of any multivariate specification to improve upon the univariate benchmark out-of-sample.

The domestic Phillips curve hypothesis receives no empirical support. Unemployment coefficients in Model 3 are uniformly insignificant, the GETS procedure eliminates all unemployment lags without exception, and the Granger causality tests in Section~\ref{sec:granger} find no predictive content from Danish labour market conditions to Danish inflation at any tested lag structure. This is not a marginal result: it holds across three distinct estimation frameworks, each approaching the hypothesis from a different angle. The finding is consistent with the broader literature documenting the flattening of the Phillips curve in small open economies \parencite{stock_watson_2008}, and has a structural explanation specific to Denmark. Under ERM\,II, monetary policy is subordinated to the exchange rate objective, which largely disables the transmission mechanism through which unemployment would ordinarily dampen inflation via interest rate adjustment. Danish inflation is instead anchored to the European price level, rendering domestic slack conditions a secondary determinant of inflation dynamics. 



This anchoring is confirmed by the GETS results. Contemporaneous EU27 inflation is the sole external regressor retained after reduction, with a coefficient of $0.716$ ($t = 8.28$, $p < 0.001$). The elimination of the surge dummy alongside the survival of the EU27 channel is instructive: the 2021--2023 inflation episode was not a distinctly Danish phenomenon requiring a separate structural break correction, but a European-wide price movement that the ERM\,II channel already captures endogenously. This interpretation rests on the assumption that Eurostat releases EU27 inflation figures earlier in the month than the Danish national release; if only the previous month's EU27 figure were available in real time, the contemporaneous channel would be infeasible and out-of-sample performance may differ. The September 2021 break identified by the QLR test in Section~\ref{sec:QLR} is therefore better understood as a large co-movement shock transmitted from the euro area than as an autonomous shift in the Danish inflation process.


The central puzzle is that EU27 inflation substantially improves in-sample fit but does not translate into out-of-sample forecast gains. The pre-surge ARIMA achieves the lowest RMSE at $0.462$, outperforming both multivariate specifications ($0.552$ and $0.546$) and the naive ARIMA ($0.554$). This disconnect is a well-documented feature of macroeconomic forecasting \parencite{stock_watson_2008}: relationships that appear stable over a historical estimation period can become redundant during regime transitions. The 2024--2026 test period represents exactly such a transition, as European inflation normalised rapidly following the post-COVID surge. The EU27 spillover coefficient, calibrated on a period of large and persistent co-movements, is better suited to tracking surge dynamics than to predicting the subsequent disinflation. The pre-surge ARIMA faces no such calibration risk: estimated on a clean sample that excludes the surge, its parameters 
reflect a long-run mean more representative of the test environment. A further caveat is that the reported RMSE averages over the full 24-month horizon: forecast errors accumulate rapidly with the horizon and all models revert toward their estimated mean at long ranges, so the ranking may partly reflect mean-reversion behaviour rather than genuine short-run predictive  content.

Several limitations qualify the conclusions. First, the 24-month test period covers a single post-surge normalisation episode with a fixed forecast origin, meaning the model ranking may not generalise to other inflation regimes and the reported metrics weight all horizons equally despite rapidly deteriorating reliability beyond the first few months. Second, the conditional forecasting approach for Models 3 and 4 assumes perfect foresight of future EU27 inflation and unemployment, overstating their practical real-time value; in a genuine out-of-sample exercise these regressors would themselves require forecasting, likely further diminishing the advantage of multivariate specifications. Third, the contemporaneous EU27 channel in Model 4 presupposes that Eurostat figures are available within the same month as the Danish release; if this timing assumption fails in practice, the channel is infeasible and the in-sample relevance of EU27 inflation may not translate into real-time forecast improvements. Fourth, the residual distributions of the Phillips Curve and GETS models display heavier tails than normality implies, as shown in Figures~\ref{fig:phillips_curve_residuals} and~\ref{fig:gets_residuals}, which may affect finite-sample inference; significance conclusions and the GETS elimination procedure should therefore be interpreted with caution, notwithstanding the large estimation sample and the use of White-robust standard errors in Model 4.